% ================================================================
%  第 4 章 引言(后半) — 原书 p.114-122, 边码 P.IX-P.XXII
%  源文: _work/src1/task-ch04b.md
%  承接 chapters/part1-ch04.tex(该文件止于 4.1.2 之末)
% ================================================================

% ---- [补译] 4.1.3 节余下部分(p.113; part1-ch04.tex 中缺此节正文) ----
\subsection{3. 罗盘与行囊}

我还该就与《收获与播种》一起构成本卷、并与它同名的另外两篇文字说几句话。

《纲领概要》勾画出我近十年来所追寻的数学思考的主要主题。我打算今后几年里,至少把其中若干稍加展开,写成一连串不拘形式的反思——我已有机会提起过它们,即“数学反思”。这份概要是去年一月我写的一份报告的逐字复制,当时是为了支持我申请国家科研中心(CNRS)的一个研究员职位。我把它收进本卷,是因为这个纲领显然远远超出我卑微个人之力所能及——哪怕上天让我再活一百年,并且我选择把这一百年用来尽可能远地推进这些主题。

《主题概要》写于1972年,是为另一次申请(法兰西学院的一个教授职位)。它按主题勾画了我当时视为自己主要数学贡献的东西。这篇文字带着它写作时心境的印记——那阵子我对数学的兴趣,说得客气些,也只能算是极边缘的。所以这份概要顶多不过是一份干巴巴、有条有理的罗列(幸好它并不求详尽……)。它看不出受到任何眼光的照耀,也看不出被任何欲望的气息吹动——仿佛我在其中像交差一样逐一过目的那些东西(而这确实就是我当时的心境),从来不曾被一种活的眼光触碰过,也从来不曾有过那种激情:要把它们带到光天化日之下,而它们当时还只是躲在雾与影的帷幕背后被隐约预感着……

然而,我之所以还是决定把这份恐怕没什么灵感的报告收在这里,主要是想堵住(假如这事办得到的话)某些高来高去的同行和某种风气的嘴:自从我离开那个我们曾共同所属的世界以来,他们总爱居高临下地看他们亲切地称之为“格罗滕迪克货色”(grothendieckeries)的东西。据说,这词的意思是死抠一些太琐碎的小事,以致一位严肃而有品味的数学家不肯为它们浪费自己当然宝贵的时间。也许这份“难以消化的文摘”在他们眼里会显得更“严肃”一些!至于我笔下那些由一种眼光、一种激情所驱动的文字,它们并不是写给那些人的——某种风气把他们维持、并且为他们开脱在一种自得自满之中,使他们对那些令我着迷的东西毫无感觉。如果说我写东西不只是为自己,那是为这样的人:他们不觉得自己的时间和自己本人珍贵到了不能无休无止地去追寻那些谁都不屑一顾的显而易见之物,不能为每一件被发现之物的内在之美而欢欣,并且把这美与我们所知的任何别的、各有其自身之美的东西区分开来。

% ---------------- p.114 ----------------
\origpage{IX} 如果我愿意把构成本卷的三篇文字彼此定位,并定出各自在这趟旅行中——我随《数学反思》一同登上了这趟旅行——所扮演的角色,我就可以说:反思—见证《收获与播种》反映并描述了我踏上这趟旅行时所怀的精神,正是这精神给了它意义。《纲领概要》描述了我的灵感之源,它们为这趟驶入未知的旅行定出一个方向——当然不是说定下一个目的地——有点像一只罗盘,或者一根结实的阿里阿德涅线。最后,《主题概要》匆匆过了一遍一份行囊,那是我1970年之前作为数学家的往昔里积攒下来的;其中至少有一部分,在旅行的这个或那个阶段会有用、会受欢迎(就像我的上同调代数与 topos 代数那套本能反应,如今在《追寻栈》里已是我不可或缺的东西)。而这三篇文字相随的次序,连同它们各自的篇幅,恰好(并非我刻意安排)反映出我在这一旅程中给予它们的分量与重视——这趟旅程的第一程,已近尾声。

\subsection{4. 一场追寻显而易见之物的旅行}

我还该就这趟一年多一点以前启程的旅行——即“数学反思”——再说几句更周详的话。关于我怀着怎样的精神踏上这趟旅行,我在《收获与播种》开头八节(即这场反思的第I、II部分)里已说得相当详细;我想,这种精神从现在起就在本卷里,也在接续它的那一卷(即《模型史》,它是《追寻栈》的第1卷)里显露出来了——后者正在完成中。因此,我觉得在这篇引言里再就此发挥,已是多余。

我当然无法预言这趟已经启程的旅行会是什么样子;这事我会随着它一路走下去而一点点发现。眼下我并没有一个哪怕只及大要的预定路线,而且我怀疑近期内会冒出这么一条来。正如我前面说过的,那些大概会启发我思考的主要主题,都多少勾画在《纲领概要》里了——那份“罗盘文本”。在这些主题当中,还有《追寻栈》的那个主要主题,即“栈”;我希望就在今年之内把它(并且就此打住)走完一遍,写成两卷,也许三卷。关于这个主题,我在《纲领概要》里写道:“……这有点像一笔债,我要向一段科学的过去偿还:在那十五年里(1955至1970年间),上同调工具的发展始终是我在代数几何基础方面工作中不变的主旋律”。因此,在预定要写的主题当中,这一个最深地扎根于我的科学“过去”。也正是这一个,在逝去的这十五年里始终作为一种遗憾在场——它也许是我离开数学舞台时留下的未竟工作中最显眼的一处空缺,而我旧日的学生或朋友里,没有一个人费心去填补它。关于这项进行中的工作,感兴趣的读者可以去查阅《纲领概要》中相应的那一节,或者去查阅《追寻栈》第一卷(正在完成中)的引言(这一回是真正的引言!)

作为我的科学过去留给我的另一份遗产,尤其还有那个我格外放在心上的 Motif 观念:它至今仍等着走出那一直被囚于其中的黑夜——尽管它现身已有足足十五年了。如果没有任何比我位置更好的人(更年轻,而且拥有他所掌握的工具与知识)决定在最近这几年里动手,那么最终由我自己来干这里所必需的基础工作,也并非不可能。

我趁此机会说明:Motif 观念的命运(或者不如说,厄运……),以及我发掘出来的另外几个观念——它们在我看来(潜在地)是其中最多产的——的命运,构成了一篇将近二十页的回顾性反思,是《收获与播种》诸“注”中
%
% ---------------- p.115 ----------------
\origpage{X} 最长的一篇(也是最后几篇之一)\fnote{3}{}{这条双注(第46、47号)及其子注已收入《收获与播种》第二部分“葬礼”,后者是它的直接续篇。}。事后我把这条注分成两部分(《我的孤儿们》与《拒绝一份遗产——或一个矛盾的代价》),此外还有紧随其后的三条“子注”\fnote{4}{}{这里指的是第48、49、50号子注(第48$'$号注是后来补加的)。}。这连续五条注合起来,是《收获与播种》里唯一不作泛泛一提、而是真正谈论数学概念的部分。这些概念在这里成了机会,用来例示数学家世界内部的某些矛盾——而这些矛盾本身又映照出人本身之中的矛盾。有一个时候我曾想把这株触手蔓生的注从它所属的文本里分出去,归入《主题概要》。那样做的好处,是能给后者一个透视,并给一篇实在太像目录的文字注入一点生气。然而我还是没有这样做,为的是保住一份见证的真实性——而这条巨型注,不管我乐意不乐意,确实就是这份见证的一部分。

\origpage{XI} 关于《收获与播种》里所说我以何种心境着手“反思”,我想在这里只补一句话——对此我已在其中一条注(《年轻人的势利——或纯洁的捍卫者们》)里表过态,那里我写道:“我毕生作为数学家的抱负,或者不如说我的欢乐与激情,始终是去发现那些显而易见之物;在眼下这部著作里,这也依然是我唯一的抱负”(《追寻栈》)。对于我一年来随“反思”进行的这趟新旅行,这也同样是我唯一的抱负。开启这趟旅行的这部《收获与播种》也不例外(至少对我的读者而言是如此,如果真有读者的话)。

\subsection{5. 一笔情愿的债}

我想以几句关于本卷《收获与播种》两则题献的话,来结束这篇引言。

那则题献——“献给曾是我的学生们的人,我曾把自己最好的给了他们——也给了最坏的”——至少去年夏天就已在我心里,尤其是在我写开头四节的时候——那四节当时还该是一本数学书的引言。这说明我很清楚(其实已经清楚好几年了):有一个“最坏的”东西需要审视——而现在正是时候,再没有比现在更好的时候了!(不过我万万没想到,这个“最坏的”最终会领着我走完一场将近两百页的冥想。)

相反,“献给曾是我的长辈们的人”这则题献,是半路上才出现的,就像这场反思的名字本身(它后来也成了一卷的名字)一样。这则题献向我揭示出他们在我作为数学家的生命中所扮演的重要角色——这角色的效力至今仍然活着。这一点,在后面的篇幅里大概会相当清楚地显现出来——所以在这里不必就此多说。这些“长辈”,按他们在我二十岁那年(大致)进入我生命的次序,是亨利·嘉当、克劳德·谢瓦莱、安德烈·韦伊、让-皮埃尔·塞尔、洛朗·施瓦茨、让·迪厄多内、罗歇·戈德芒、让·德尔萨特。当时那个初来乍到的我,受到了他们每一个人的善意接待;此后他们中许多人给了我持久的友谊与情谊。我还须在此提到让·勒雷:在我与“数学家世界”的第一次接触时(1948/49年),他的善意接待同样是一种珍贵的鼓励。我的反思显现出一笔感恩之债,欠他们每一个人——这些“来自另一个世界、另有另一份命运”的人。这笔债绝不是负担。发现它时,它像一份欢乐来临,让我变得更轻。

\noindent 1984年3月底\par

% ---------------- p.116 ----------------
\section{(II) 一桩敬重之举}

(5月4日—……6月)

\subsection{6. 葬礼}

\origpage{XII} 一件意料之外的事,把一场本已走到头的反思又重新发动起来。它开启了一连串大大小小的发现,在过去这几周里逐步揭开了一个一直模糊不清的局面,并使它的轮廓愈发分明。这尤其使我周详而深入地走进了一些事件与处境——此前它们只是被顺带一提,或隐约暗示过。于是,前面提到过的那篇关于一部作品之沉浮的“十五页回顾性反思”(引言,第4节),就获得了意想不到的规模,多出了大约两百页。

由于事势所迫,也由于一场反思的内在逻辑,我在路上被引到把别人牵涉进来,其程度不亚于牵涉我自己。被牵涉得比谁都深的(除了我自己),是一个与我有着将近二十年友谊的人。关于他,我曾写道(那是一种委婉说法\fnote{5}{}{关于这个“委婉说法”的含义,见注《与众不同者》,第57$'$号。}):在这份扎根于共同激情的亲切友谊的最初几年里,他“多少有点像个学生”;而在很长一段时间里,在我内心深处,我把他看作某种“合法继承人”——所继承的,是我相信自己还能在数学上带来的东西,是在那部始终零碎的已发表作品之外的东西。很多人早已认出他是谁了:他就是皮埃尔·德利涅。

我并不为随这些注一起公开一份对一段私人关系的个人反思、并且就这样未经征求他同意而把他牵涉进来而道歉。一个长期隐而不显、混乱不清的局面,终于被摊到光天化日之下加以审视,在我看来是重要的,对所有人都是健康的。这样做,我是在提供一份见证——无疑是主观的,既不自称穷尽了一个微妙而复杂的局面,也不自称没有错误。它的首要长处(正如我过去那些出版物、或眼下正在写的那些的首要长处一样)是它存在着,供可能对它感兴趣的人取用。我所操心的,既不是要说服人,也不是要躲在那些所谓“确凿”的东西背后,以免于错误或怀疑。我所操心的是要真实——把我在每一刻所看见、所感到的东西就那样说出来——以此作为深入它们、理解它们的手段。

“葬礼”这个名字——用于所有那些与“一段过去的重量”相关的注的总称——在反思过程中以一种越来越强的力量自显出来\fnote{6}{}{在这场反思接近尾声时,另一个名字浮现出来,表达了某幅图景的另一个触目惊心的方面——那幅图景在过去五周里逐步在我眼前展开。那是一个童话的名字,我将在相应之处回到它:《中国皇帝的袍子》……}。在其中,我扮演那位预先的死者的角色,与几位(年轻得多的)数学界同行组成凄凉的送葬队伍;他们的作品落在我1970年“离去”之后,并且凭着某种风格、某种对数学的进路,带着我影响的印记。这些人里排在头一个的,是我的朋友左格曼·梅布胡特;他有幸——这运气沉重得很——要承受一个被当作“1970年之后的格罗滕迪克的学生”的人所遭遇的一切不利,却又没享受到与我接触、得到我鼓励与指点的那种好处,因为他只是通过我的文字、从我的作品那里“受教”的。那阵子(在他所出没的那个世界里),我已经俨然一副“死者”的模样,以致长期以来,连见上一面的念头显然都不曾出现过;一段持续的往来(无论私人的还是数学的),直到去年才终于结上。

% ---------------- p.117 ----------------
\origpage{XIII} 这并没有妨碍梅布胡特逆着一种专横的风气、逆着他那些(曾是我的学生的)长辈们的轻蔑,在近乎完全的孤立中,做出独创而深刻的工作——靠的是一种意想不到的综合:把佐藤学派的想法与我的想法合到一起。这项工作提供了一种新的入手之处,可以把握解析簇与代数簇的上同调,并预示着我们对这种上同调的理解将有一场大规模的更新。毫无疑问,假如梅布胡特在那些最适合做这件事的人那里,得到了他们当年在我这里得到过的热情接待与毫无保留的支持,那么这场更新早已是既成之事,而且已经完成多年了。至少,自1980年10月起,他的想法与工作为代数簇上同调理论的一场引人注目的重启提供了灵感与技术手段,使这门理论终于(除了德利涅围绕韦伊猜想的那些结果之外)走出了一段漫长的停滞期。

\origpage{XIV} 难以置信却又千真万确:近四年来,他的想法与结果被“所有人”使用着(与我的想法同等待遇),而他的名字仍被小心地无视、被默不作声地略过——恰恰是被那些第一手了解他的工作、并在自己的著作里根本性地使用它的人。我不知道数学界在别的什么时候可曾见过这样的失格:其中某些最有影响、最有声望的成员,在普遍的漠然中,带头蔑视数学家这行伦理中最被普遍承认的规矩。

我看见四个男人,才能出众的数学家,他们和我一起享有——也有资格享有——这场以沉默与轻蔑举行的葬礼的哀荣。而我在每个人身上都看见:轻蔑如何啃咬着那曾激励过他的美好激情。

除了这几个人,我尤其看见两个人:他们一个接一个站在数学公共广场的聚光灯下,在众人的簇拥中主持葬礼,而同时(在一种更隐蔽的意义上)他们也被埋葬——被自己的双手埋葬,与他们刻意埋葬的那些人一同被埋葬。其中一个我已经点过名了。另一个同样是从前的学生、从前的朋友:让-路易·维尔迪耶。1970年我“离去”之后,他与我之间的接触没有维持下来,除了职业层面上几次匆促的会面。所以大概正因如此,他在这次反思里只通过他职业生涯中的某些举动出现;至于这些举动可能的动机——就他与我的关系这一层面而言——则未经审视,而且我对此也完全摸不着头脑。

如果说有一个追问在逝去的这些年里始终逼着我,它曾是《收获与播种》的一个深刻动机,也在这整场反思中一路跟着我,那就是:在某种精神、某种风气的到来之中,我占了怎样一份——正是这种精神与风气,使得像我说的那样的失格成为可能,而这一切发生在一个曾属于我的世界里,在我作为数学家的二十多年里,我曾与它认同为一。这场反思让我发现:由于我身上某些虚浮的态度——它们表现为对才具平平的同事一种无言的轻蔑,表现为对我自己、对某些才具出众的数学家的自我迁就——我对这种精神并不陌生;而今天,我看见这精神在我曾爱过的那些人当中肆意铺展,也在我曾把自己热爱的一门手艺教给的那些人当中铺展;那些我没有好好爱、也没有好好教的人,如今在这个曾为我所珍爱、而我已经离开的世界里定调子(当他们不是在立法的时候)。

我感到一阵风刮起来,那是自得、玩世不恭与轻蔑的风。“它刮着,不管什么‘功绩’什么‘过失’,用它的气息烧灼卑微的志向,也烧灼最美的激情……”。我明白了:这阵风,是我参与播下的盲目而不经心的种子那丰饶的收成。而如果它的气息回吹到我身上,回吹到我曾交托给别人手里的东西上,回吹到我今天所爱的人身上,而且
%
% ---------------- p.118 ----------------
\origpage{XV} 回吹到那些敢于自称属于我、或仅仅是受我启发的人身上,那么这是事物的回还,我没有理由抱怨,而且它有许多东西要教我。

\subsection{7. 葬仪的序次}

于是,在“葬礼”这个名目下,我在目录里把那些与那个表面平淡无奇的节“一段过去的重量”(第50节)相关的主要“注”的浩浩荡荡的队列归拢在一起,从而让这个名字充分显出了它的意义——当初写《收获与播种》“初稿”的这最后一节时,这名字一上来就落到了我头上。

\origpage{XVI} 在这支亲缘关系错综的长长的注的队列里,过去四周里才加入的那些(注(51)至注(97))\fnote{8}{}{还须加上1984年5月12日的第104号注。第98号及其后的各注(上一句提到的第104号注除外)构成反思的“第三口气”,自1984年9月22日起;它们也标了日期。}显得与众不同:它们是唯一标了日期的(4月19日至5月24日)\fnote{9}{}{在同一天写的一连串连续的注里,只有第一条标日期。其余未标日期的注是第44$'$号至第50号(构成第I、II、III支队伍)。第46、47、50号注写于3月30日或31日,第44$'$、48、48$'$、49号注写于四月上半月,最后,第44$''$号注是标了日期的(5月10日)。}。我觉得最自然的做法,是按它们在反思中相继出现的时序\fnote{10}{}{有时我在这个时间顺序里作过幅度不大的调换,换成一个“所谓逻辑的”次序,只要我觉得反思进程的整体印象不会因此失真。不过,作为仅有的例外,我要指出十一条注(其编号前有一个感叹号!)——它们源自某条注之后的页底脚注,篇幅涨到了让人却步的规模,于是我把它们各自放在与之相关的那条注之后(除了第98号注,它关联的是第47号)。}把它们排出来,而不是按某种所谓的“逻辑”次序,也不是按先前各注中提到它们的先后次序。为了能找回后一种(绝非线性的)承继关系——即各条参与其中的注之间的谱系——我在目录里让每条的编号后面跟着(在它之前的那些注中)第一次提到它的那一条的编号\fnote{11}{}{当对某条注的引用(如(45))就出现在“一段过去的重量”这一节本身之中时,则把这一节的编号(50)加括号,置于该注编号之后,如“46(50)”。},若没有这种情况,则跟着它所直接续写的那一条的编号\fnote{12}{}{在目录里,若某条注是前一条注的直接续写(此时两号相连),其编号前加一个 * 号。例如“*47, 46”表示第47号注是第46号注的直接续写(而这里第46号并非紧邻在它前面的那一条,紧邻的是第46\textsubscript{9}号注)。\par 最后,我在目录里给那些后面不跟别的编号的注加了着重号,也就是说,给那些代表反思“新起点”的注——它们并不嵌在已经做过的反思的某个确定位置上。}。(后一种关系在正文里另有标示:在前一条注的末尾放一个指号,如注(46)末行末尾的($\Rightarrow$ 47),指向续写它的注(47)。)最后,某一注中某些稍具技术性的细节,则集中在该注末尾,编成子注,用接在该注原编号之后的连续下标标示——如注(46)《我的孤儿们》的子注(46\textsubscript{1})至(46\textsubscript{9})。

为了多少给“葬礼”的整体安排立起结构,也为了让人在蜂拥而至的众多注里能认得出路,我觉得在此场合宜于在队列中插入几个庄重而富于暗示的小标题,每个小标题都走在并引领一支或长或短的送葬队——由主题相同、彼此相连的连续若干注组成。

于是我有幸看见,在一支前来为我送葬的庄严长队里,十支\fnote{13}{}{(9月29日)其实最终有十二支队伍,若把灵车(X)和“死者(仍未过世)”(XI)算进去的话——后者在最后关头又溜进了队列……}队伍一支接一支地聚拢起来——有的卑微,有的气派,有的悔恨,有的暗暗欢腾,在这样场合下本也免不了如此。于是走上前来的有:身后的学生(人人都把无视他当作本分)、孤儿们(为这场合才刚被掘出)、时风及其名流(这是我活该)、诸Motif(我所有孤儿中最晚出生、也最晚被掘出的),
%
% ---------------- p.119 ----------------
\origpage{XVII} ……我的朋友皮埃尔谦逊地领着最要紧的那一支队伍,紧随其后的是“一致同意”——那些(默默地)协同一致的注——以及全体出动的(所谓)“乖戾”研讨会(它与身后的学生、即“无名学生”划清界限的方式,是夹在中间的几支捧着鲜花与花环的送葬队);最后,为了给这气派的队列一个体面的收场,走上前来的还有“学生”(绝非身后的,更绝非无名的),又名“老板”,身后跟着我那一帮忙忙叨叨的学生(带着大量铁锹与绳索),最后是灵车(车上赫然四口漂亮的橡木棺材,用螺丝牢牢拧紧,还不算那个掘墓人)……十支队伍终于齐了(总算齐了),缓缓走向葬礼仪式。

\origpage{XVIII} 仪式的高潮是悼词,由我的朋友皮埃尔亲自出马,以完美的分寸献上——他应众人之请主持葬礼,人人满意。仪式在一曲终局而决定性的“自深处”(De Profundis)歌声中结束(至少人们这样希望)——这歌由可敬的死者本人唱出,作为真诚的感恩;他瞒着所有人,从自己那场声势骇人的葬礼上活了下来,并且按说从中得了教训,心满意足——这份满足,便构成了那场值得纪念的葬礼的最后一个音符、最后一个和弦。

\subsection{8. 一个秘密的终结}

在这场(但愿是)最后的反思阶段中,我看出有必要在《数学反思》第1卷的“附录”里,除前面谈到的三篇之外,再附上另外两篇性质属于数学的文字\fnote{14}{}{此外,我想给《主题概要》(见引言第3节“罗盘与行囊”)加一份“评注”,就我对其中粗略过目的那些“主题”的贡献作些说明,也就那些在我的数学作品中主要思想力的生成中起过作用的影响作些说明。最近六周的回顾已经显出(令我自己也吃惊):在大多数这些思想的启动上,以及在我1955至1970年间给自己定下的某些“大任务”上,塞尔都起了“雷管”的作用。\par 最后,还有一篇性质属于(通常意义上的)数学的文字,也是唯一(偶然)出现在非技术性文本《收获与播种》里的——我要指出注《屠杀》(第87号)下的子注第87\textsubscript{1}号:我在那里以它应有的仔细,阐明凝聚层语境中为人熟悉的黎曼–罗赫–格罗滕迪克定理的一个“谨慎的”(猜想性)变体。这个猜想(与另外若干一起)出现在1965/66年度 SGA~5 讨论班的收尾报告里;那份报告(如同许多别的报告一样)在十一年后以 SGA~5 之名出版的文集中没有留下任何痕迹。这个关键讨论班落在我某些学生手里之后的种种遭际,以及这些遭际与某种“SGA~4½ 行动”的关联,在注第63$''$、67、67$'$、68、68$'$、84、85、85$'$、86、87、88号所进行的反思过程中逐步显露出来。\par 又及:另一条给出相当充实的数学评注的注,讨论的是为那些已知的、拥有所谓“六种运算”对偶形式的情形(在可能的范围内)抽出一个共同“topos 式”框架是否合宜——我还要指出注《赊账的论文与全险》(第81号)下的子注第81\textsubscript{2}号。}。

第一篇是一份带评注、分两部分的报告的复制,是我1968与1969年就 P.\ 德利涅的工作(其中有些至今仍未发表)写的,对应于1965/67/68三年间在法国高等科学研究所(IHÉS)的数学活动。

另一篇是一份“六种变率公式表”的草稿,汇集了我在1956至1963年间抽出的某种对偶形式(灵感来自庞加莱对偶与塞尔对偶)的共同特征;这份公式表后来表明,对于迄今遇到的一切上同调对偶情形,它都具有“普适”的品格。随着我离开数学舞台,这套形式似乎就废弃不用了,以致据我所知,还没有人(除我之外)肯费神哪怕只写出那些基本运算的清单、它们所给出的基本典范同构的清单,以及这些同构之间那些本质的相容关系的清单。

这份融贯的公式表草稿,对我将是显然的第一步,通向那幅“Motif 之梦的宏伟全景图”——十五年来,它一直“等着那位敢于挥笔把它画出来的数学家”。

% ---------------- p.120 ----------------
\origpage{XIX} 看样子,这位数学家不是别人,正是我自己。确实,该是时候了:将近二十年前诞生、在私密中交托出去的东西——它本不是为了成为一个人的特权,而是为了供所有人取用——终于要走出秘密的黑夜,再一次诞生于白昼的充足光线下。

确实,除我之外,只有一个人对这套“Motif 的瑜伽”有贴身的了解——在我离去之前的那些日子里、那些年里,他一天天、一年年从我口中把它学了去。在我有幸发现并带到光天化日之下的一切数学事物当中,这个 Motif 的实在在我眼里仍是最迷人的、最充满奥秘的——它就在“几何”与“算术”之间那深层同一性的核心处。而这套“Motif 的瑜伽”——那长久被忽视的实在把我引向的瑜伽——也许是我在作为数学家的生命头一段里抽出的最有力的发现工具。

但同样真实的是:这个实在,以及这套努力把它圈得尽可能近的“瑜伽”,从来不曾被我守为秘密。我被撰写基础这一必须完成的任务占住了(此后人人都很乐意把它们原封不动地用在自己日常的工作里),没有拿出那必不可少的几个月,为这套 Motif 的瑜伽写一份宏阔的整体草稿,从而把它交到所有人手里。然而在我出乎意料地离去之前的那些年里,我并没有放过机会——每逢相遇,就向愿意听的人讲它,首先是我的学生们;而他们(其中一人除外)都把它忘了,就像所有人都把它忘了一样。我讲它,不是为了推销什么会挂上我名字的“发明”,而是为了引人注意一个每走一步都显现出来的实在——只要你关心代数簇的上同调,尤其是关心它们那些“算术”性质,以及迄今已知的各种上同调理论彼此之间的关系。这个实在,与当年“无穷小”的实在一样可触——那东西被感知到,远早于严格语言的出现,而正是那种语言使人能完满地把握它并把它“确立”下来。而为了把握 Motif 的实在,我们今天绝不缺一套灵活而合用的语言,也不缺在数学理论建构方面成熟老练的经验——这些,我们的前人是缺的。

如果我从前在屋顶上大声喊过的东西落进了聋耳朵,如果某一个人的沉默不语所收到的回声,是所有那些装作对上同调感兴趣的人(而他们明明和我一样有眼睛有手……)的沉默与麻木,那么我也不能把责任只归到那一个选择把“好处”留给自己的人身上——那本是我为所有人而交托给他的。我们不得不承认:我们这个时代,其科学产出的狂飙足以与投入军备或消费品的劲头相比,却离我们十七世纪前人那种“大胆的活力”相去甚远——他们为了发展无穷小的计算,“不走那么多弯路”,不会因为这东西是否“猜想性”的顾虑而停住,也不等他们中某位显赫人物肯来给他们开绿灯,才去抓住每个人自己亲眼看见、第一手感到的东西。

\subsection{9. 场景与角色}

凭它自身的内部结构,也凭它那个特殊的主题,“葬礼”(如今已占《收获与播种》全文一半以上)在很大程度上、并且从逻辑观点看,独立于它前面那场漫长的反思。然而这只是一种纯粹表面的独立。对我来说,这场围绕着一场“葬礼”的反思——那葬礼渐渐从“未言”与“预感”的雾中走出来——与在它之前、它由之而生并且由之获得全部意义的那场反思,是不可分的。它起头时只是“顺路”匆匆一瞥一部作品的沉浮——那部作品我有点(很大程度上)已经看不见
% ---------------- p.121 ----------------
\origpage{XX} 了;后来,既未预谋也未寻求,它变成了一场对一段重要关系的冥想,而这又把我引向对那部作品在“曾是我的学生们”手中的命运的反思。把这场反思从它自发从中生出的那场反思里割出去,在我看来就是把它降格为一幅简单的“风俗画”(甚至是在数学“上流社会”里算的一笔账)。

\origpage{XXI} 确实,若有人偏要如此,整个《收获与播种》也能作同样的降格,变成一幅“风俗画”。当然,某个时代、某个圈子里通行的风俗,那些参与塑造该圈中人生活的风俗,自有其重要性,也值得描写。然而《收获与播种》的细心读者会明白:我的打算不是描写风俗,也就是说,不是描写某一个场景——它随时代而变、因地而变,我们的行动就在它上面展开。这场景在很大程度上界定并划出了我们身上各种力量可用的手段,使它们得以表达。场景以及它提供的手段(还有它强加于人的“游戏规则”)千变万化,而我们身上那些深层力量的本性——它们在集体层面上塑造场景,在个人层面上又在场景上表达自己——看来从一个圈子、一种文化到另一个圈子、另一种文化,从一个时代到另一个时代,都是同一个。如果说我这一生里,在数学之外、在爱女人之外,还有什么东西让我感到奥秘与吸引(不错,是晚近的事),那正是这少数几种力量的隐秘本性——它们有能力驱使我们行动,为“最好”也为“最坏”,为埋藏也为创造。

\subsection{10. 一桩敬重之举}

这场最终取名为“葬礼”的反思,起初是一桩敬重之举。敬重那些我发现了的东西——我看着它们在一片虚无中凝聚、成形,我是第一个尝到它们的滋味与劲道的人,我给它们起了名字,既为表达我对它们的认识,也为表达我的敬重。对这些东西,我给出了自己身上最好的部分。它们汲取着我身上的力量,生长、绽放,像许多粗壮的枝条从同一根活着的树干上迸发出来,那树干有着粗壮而众多的根。这些是活着的、在场的东西,不是可有可无的发明——它们紧密相依,构成一个活生生的统一体,这统一体由它们每一个构成,又给每一个以位置与意义、一个起源与一个终点。我很久以前就把它们留下了,没有丝毫不安,也没有丝毫遗憾,因为我知道我留下的东西是健全而强壮的,它根本不需要我,就能依着它自己的本性继续生长、绽放、繁衍。我留下的不是一袋金币——那是可以被人偷走的;也不是一堆工具——那是会生锈、会腐烂的。

然而,年复一年,当我自以为离那个我离开的世界很远的时候,却有一阵阵阴险的轻蔑、隐隐的嘲弄,一直飘到我隐居之处,在这里那里回旋,指着这些东西中的某些——我明知它们强壮而美,有它们的位置,有它们独一无二、任何别的东西都永远无法代替的功能。我感到它们像是敌对世界里的孤儿——那个世界患着轻蔑的病,专门纠缠没有铠甲护身的东西。这场反思就是在这样的心境中开始的,作为对这些东西的一桩敬重之举,并因此也是对我自己的一桩敬重之举——作为对它们与我之间那条深沉纽带的提醒:谁若乐于对这样一件由我的爱滋养出来的东西摆出轻蔑的样子,他所乐于轻蔑的正是我,以及一切从我这里生出的东西。

对于另一种人,情形也一样:他第一手知道那条把我和某件东西连起来的纽带——那东西他不是从我这里、又是从哪里学来的呢——却摆出一副把这条纽带视为无足轻重、或者干脆无视它的样子,或者(哪怕只是默默地、以省略的方式)为自己、或为别人僭取一种虚假的“父权”。我在其中
% ---------------- p.122 ----------------
\origpage{XXII} 清清楚楚看到一桩轻蔑之举:轻蔑那件从工匠手里生出的东西,轻蔑那默默而精细、使这件东西得以生出并由此诞生的工作,轻蔑那工匠,而首先(以一种更隐蔽也更本质的方式)轻蔑他自己。

如果我的“回归数学”只消让我重新想起这条纽带,只消在我心里唤起这桩当着所有人——当着那些摆出轻蔑样子的人,也当着无动于衷的见证者——的敬重之举,那么这次回归也就不算白费。

确实,我与我留下的那部写出来的和没写出来的(至少是未发表的)作品,真的失去了接触。开始这场反思时——我看那些枝条还相当分明,却不怎么记得它们原是同一棵树的一部分。说来奇怪:非得那幅我留下之物遭劫掠的图景在我眼前一点点展开,我才在自己心里重新找回那种意义——那种关于被如此劫掠、如此拆散之物的活生生的统一性的意义。一个人拿走了一些金币,另一个拿走了一两件工具,用来炫耀,甚至用来使用——但是那赋予我所留之物以生命和真正力量的那种统一性,却从每一个人、也从所有人手里漏掉了。不过我清楚知道有一个人,深深感到过这种统一性和这种力量,在他内心深处今天仍然感到它;而他却乐于耗费自己身上的力量,要去摧毁他在别人身上、通过别人的作品感到过的那种统一性。作品的美与创造之德,恰恰就在这种活生生的统一性之中。尽管遭了劫掠,我仍把它们找回来,完好如初,仿佛我刚离开它们——只是我成熟了,今天用一双新的眼睛看它们。

然而,如果说有什么东西真被劫掠、被损毁、被卸掉了它原本的力量,那是在那些人身上:他们忘了自己身上安歇着的力量,自以为在劫掠一件任由他们摆布的东西,其实他们只是把自己与那创造之德割断——那东西就在他们手边,一如它在所有人手边,却绝不任他们摆布,也不在任何人的权力之下。

这样,这场反思,以及通过它而来的那次不期而至的“回归”,也让我重新接触到一种被遗忘的美。正是完完全全感受到了这种美,才使那桩敬重之举有了全部意义——它在注《我的孤儿们》\fnote{15}{}{这条注(第46号)在收入“葬礼”的所有注中,按时间是最早的一条。}里表达得笨拙,而我刚才在这里,在完全知情的情况下,又把它重述了一遍。
